\chapter{高频、微观结构与执行分层答案}
\begin{enumerate}[leftmargin=*]
  \item 日历时间连续推进，事件时间按离散消息计数。条件强度只允许使用事件发生前的滤过 $\mathcal F_{t^-}$。
  \item $\alpha/\beta<1$ 保证平均后代数低于 1；开收盘确定性高峰若未去除，会制造虚假的自激相关。
  \item 一阶条件 $n/\lambda-T=0$ 给出 $\widehat\lambda=n/T$；本章 $n=T=4$，故强度为 1，对数似然为 $-4$。
  \item $V_t(x)=\min_q\{\eta q^2+\gamma(x-q)^2+V_{t+1}(x-q)\}$，并限制交易量与流动性；$V_T(0)=0$，其他终态为无穷。
  \item 对每个价位保存 FIFO 队列；第三张同价订单只能在前两张之后成交。部分撤单不改变其剩余数量的原始时间优先级。
  \item 在零库存风险时，凸冲击成本由均匀序列 $(2,2,2)$ 最小化，成本 12；取 $\gamma=0.3$ 后枚举得到 $(3,2,1)$，目标 17。
  \item 到达模型用似然与残差诊断；状态机用固定事件序列；策略用同一事件流上的独立基线和成本后指标。三者不能共享一个模糊失败标签。
  \item 报告行情、决策、网络与撮合延迟，价差、冲击和机会成本，库存峰值与终值，并声明隐藏单、队列位置与交易所规则未覆盖之处。
\end{enumerate}

\section*{独立 oracle 核对表}
四个间隔的强度为 1，对数似然为 $-4$。手工回放最终买一剩 1、卖一剩 2，累计成交 11。TWAP $(2,2,2)$ 与库存惩罚解 $(3,2,1)$ 的冲击成本分别为 12 与 14；在同一库存惩罚下总目标分别为 18 与 17，最优库存路径为 $(6,3,1,0)$。
